\section{\tempo{} Overview}
\tempo{} is a lightweight synchronous runtime for OCaml. It adopts the delayed
absence principle of FairThreads and ReactiveML, but implements it entirely at
library level using OCaml 5.3 algebraic effects. Programs execute over a
sequence of logical instants, interact via valued signals, and compose
concurrently under a deterministic scheduler.

The runtime preserves the core control semantics: signal presence is decided
within the instant, absence is observed only at instant end, and preemption is
weak and instant-aligned. The key design choice is to encode these constructs
as effectful operations interpreted by handlers. This keeps the model explicit
in user code and avoids changes to the language runtime, while still enabling
expressive control operators and a clear mapping to synchronous concepts.

Compared to ReactiveML, \tempo{} trades language-level syntax and compilation
support for a pure library approach based on OCaml 5.3 effects, making it
easier to embed into existing codebases and to experiment with scheduling.

\paragraph{Programming model}
\tempo{} programs are written as ordinary OCaml functions that receive two
signals: an input signal provided by the host and an output signal used to
publish results at the end of each instant. Instants are advanced by the
\texttt{execute} loop, which invokes host hooks before and after each instant.
Signals are first-class values, so they can be stored in data structures,
returned from functions, and passed between components.

\paragraph{Determinism and control}
Determinism is defined at instant boundaries: given the same input emissions,
the observable outputs at each instant are fixed, regardless of internal task
interleavings. Within an instant, the scheduler runs tasks in a fixed order,
but the semantics does not expose this ordering, which keeps the model
declarative from the user's perspective.

\paragraph{Design rationale}
\tempo{} adopts algebraic effects for three reasons. First, effects provide
first-class, delimited continuations, which allow the scheduler to control
resumption points precisely. Second, a library-level approach avoids modifying
the OCaml compiler and keeps integration friction low. Third, handlers make the
control flow explicit in user space, which simplifies experimentation with
scheduling policies and debugging.

\paragraph{Non-goals (current scope)}
\tempo{} does not attempt to provide static guarantees (e.g., single-emission
checks or determinism proofs) and does not perform compilation into automata.
These choices keep the library small and explainable, but they also motivate
the future work discussed in Section~\ref{sec:future-work}.
